\documentclass[12pt, a4paper]{report}
\usepackage[utf8]{inputenc}
\usepackage[italian]{babel}
\usepackage{amsmath, amssymb, amsfonts}
\usepackage{geometry}
\geometry{margin=2.5cm}
\usepackage{hyperref}
\usepackage{booktabs}

\title{\Huge \textbf{Elettronica Applicata}\\ \Large Appunti Integrati del Corso}
\author{\textbf{Corso di Laurea Magistrale in Tecnologie Avanzate}}
\date{\today}

\begin{document}

\maketitle

\cleardoublepage
\tableofcontents
\cleardoublepage

% ==========================================
% PARTE I: RIPASSO E FONDAMENTI DI CIRCUITI
% ==========================================
\part{Ripasso di Elettronica e Teoria dei Circuiti}

\chapter{Dai Campi Elettromagnetici alle Leggi Circuitali (QSA)}

\section{Introduzione al Corso e Convenzioni}
\subsection{Analisi vs. Sintesi Circuitale}
In ambito elettronico si distinguono due classi fondamentali di problemi:
\begin{itemize}
    \item \textbf{Analisi}: Dato un circuito topologicamente definito, determinare le correnti nei rami $I_k$ e le tensioni ai nodi $V_k$ applicando le leggi circuitali.
    \item \textbf{Sintesi (Design)}: Progettare l'architettura interna di un blocco funzionale affinché soddisfi specifiche funzioni di ingresso/uscita.
\end{itemize}

\subsection{Notazioni Formali}
Nel corso si adotterà la seguente convenzione di notazione:
\begin{itemize}
    \item $v(t), i(t)$: Grandezze tempo-varianti nel dominio del tempo (regimi alternati o transitori).
    \item $V, I$: Valori stazionari in corrente continua (DC).
    \item $V(s), I(s)$: Variabili trasformate nel dominio della frequenza complessa di Laplace $s$.
\end{itemize}

\section{Dalle Eq. di Maxwell alle Variabili Circuitali}
\subsection{Tensione, Forza Elettromotrice (FEM) e Potenziale (DDP)}
Data una carica $q$ soggetta a una forza elettrica generica $\vec{F}$, si definisce il \textbf{campo elettromotore} come $\vec{E}_m = \vec{F}/q$. Il lavoro compiuto per spostare la carica da $A$ a $B$ definisce la \textbf{tensione} $V_{AB}$:
\begin{equation}
    V_{AB} = \int_{A}^{B} \vec{E}_m \cdot d\vec{s} \quad [\mathrm{V}]
\end{equation}
La tensione misura sempre una quantità scalare tra due punti dello spazio. L'unità di misura del campo elettrico nei circuiti è spesso in $[\mathrm{V/m}]$.


Lungo un percorso chiuso, la circuitazione del campo elettromotore definisce la \textbf{Forza Elettromotrice (FEM)}:
\begin{equation}
    \xi = \oint \vec{E}_m \cdot d\vec{s}
\end{equation}

Nel solo caso di campo elettrostatico conservativo ($\vec{E} = -\vec{\nabla}\phi$), il lavoro non dipende dal percorso e la tensione coincide con la \textbf{Differenza di Potenziale (DDP)}:
\begin{equation}
    V_{AB} = \phi_A - \phi_B \implies \oint \vec{E} \cdot d\vec{s} = 0
\end{equation}

Nel caso generale non statico, la presenza di campi magnetici variabili nel tempo richiede l'uso del potenziale vettore $\vec{A}$:
\begin{equation}
    \vec{E} = -\vec{\nabla}\phi - \frac{\partial \vec{A}}{\partial t}
\end{equation}

\subsection{Densità di Corrente e Corrente Elettrica}
A livello microscopico, la densità di corrente è legata al campo elettrico dalla relazione $\vec{j} = \sigma \vec{E} = n q \mu \vec{E}$, dove $\sigma$ è la conducibilità, $n$ la densità volumica dei portatori e $\mu$ la mobilità. La corrente elettrica macroscopica $i(t)$ rappresenta il flusso di $\vec{j}$ attraverso una sezione trasversale $A$:
\begin{equation}
    i(t) = \frac{dq}{dt} = \iint_{A} \vec{j} \cdot d\vec{A} \quad [\mathrm{A}]
\end{equation}

\section{Validità delle Leggi di Kirchhoff e Condizioni QSA}
In ambito elettronico, la risoluzione dei circuiti mediante le variabili di tensione ($V$) e corrente ($I$) si basa sulle leggi di Kirchhoff alle maglie (\textbf{KVL}) e ai nodi (\textbf{KCL}). Tuttavia, tali leggi non costituiscono principi primari dell'elettromagnetismo, ma sono approssimazioni della teoria dei campi descritta dalle equazioni di Maxwell. È quindi fondamentale stabilire i \textbf{limiti di validità} entro cui è lecito adottare tali leggi.


\subsection{Limiti di Validità della KVL e Prima Condizione QSA}
Partiamo dalla III equazione di Maxwell (Legge di Faraday-Neumann-Lenz) nella sua forma locale:
\begin{equation}
    \vec{\nabla} \times \vec{E} = -\frac{\partial \vec{B}}{\partial t}
\end{equation}
Calcolando l'integrale di linea del campo elettrico $\vec{E}$ lungo un percorso chiuso $C$ (che rappresenta una maglia del circuito avente superficie $S$):
\begin{equation}
    \oint_{C} \vec{E} \cdot d\vec{s} = -\iint_{S} \frac{\partial \vec{B}}{\partial t} \cdot d\vec{A} = -\frac{d\Phi(\vec{B})}{dt} = \xi_i
\end{equation}
D'altra parte, la \textbf{Legge di Kirchhoff alle Maglie (KVL)} impone che la somma algebrica delle cadute di potenziale lungo una maglia sia identicamente nulla:
\begin{equation}
    \sum_{k} V_k = 0 \iff \oint_{C} \vec{E} \cdot d\vec{s} = 0
\end{equation}
Uguagliando le due espressioni, notiamo che la KVL risulta matematicamente valida \textbf{se e solo se} il termine di variazione temporale del flusso magnetico concatenato alla maglia è trascurabile. 

Da questa necessità deduttiva ricaviamo la \textbf{Prima Condizione di Quasi-Staticità (QSA \#1)}:
\begin{equation}
    \frac{\partial \vec{B}}{\partial t} \approx 0 \implies \vec{\nabla} \times \vec{E} \approx 0 \implies \oint_{C} \vec{E} \cdot d\vec{s} \approx 0
\end{equation}
\textit{Significato fisico}: Il campo elettrico viene approssimato come conservativo ($\vec{E} \approx -\vec{\nabla}\phi$). Si assume che il circuito non emetta né assorba radiazione elettromagnetica e che le induttanze parassite distribuite siano trascurabili, isolando la variazione di flusso magnetico unicamente all'interno degli \textbf{induttori concentrati}.

\subsection{Limiti di Validità della KCL e Seconda Condizione QSA}
Consideriamo la IV equazione di Maxwell (Legge di Ampère-Maxwell):
\begin{equation}
    \vec{\nabla} \times \vec{B} = \mu_0 \left( \vec{j} + \varepsilon_0 \frac{\partial \vec{E}}{\partial t} \right)
\end{equation}
Applicando l'operatore divergenza ad entrambi i membri ed utilizzando l'identità vettoriale per cui la divergenza di un rotore è nulla ($\vec{\nabla} \cdot (\vec{\nabla} \times \vec{B}) = 0$), otteniamo:
\begin{equation}
    \vec{\nabla} \cdot \left( \vec{j} + \varepsilon_0 \frac{\partial \vec{E}}{\partial t} \right) = 0 \implies \vec{\nabla} \cdot \left( \vec{j} + \vec{j}_D \right) = 0
\end{equation}
dove $\vec{j}_D = \varepsilon_0 \frac{\partial \vec{E}}{\partial t}$ rappresenta la \textbf{densità di corrente di spostamento}.

La \textbf{Legge di Kirchhoff ai Nodi (KCL)} richiede che la somma delle correnti di conduzione uscenti da un nodo (superficie chiusa attorno al nodo) sia nulla:
\begin{equation}
    \sum_{k} I_k = 0 \iff \vec{\nabla} \cdot \vec{j} \approx 0
\end{equation}
Confrontando le relazioni, la KCL risulta valida \textbf{se e solo se} la corrente di spostamento $\vec{j}_D$ all'esterno dei componenti è trascurabile rispetto alla corrente di conduzione $\vec{j}$.

Da questa necessità deduttiva ricaviamo la \textbf{Seconda Condizione di Quasi-Staticità (QSA \#2)}:
\begin{equation}
    |\vec{j}| \gg \varepsilon_0 \left| \frac{\partial \vec{E}}{\partial t} \right| \implies \vec{j}_D \approx 0 \implies \vec{\nabla} \cdot \vec{j} \approx 0
\end{equation}
\textit{Significato fisico}: La corrente di conduzione si comporta come un campo solenoidale (assenza di accumulo locale di carica nei nodi o nei rami), confinando l'effetto della corrente di spostamento unicamente all'interno delle armature dei \textbf{condensatori concentrati}.

\section{Limiti Quantitativi della QSA, Parametri Concentrati e Scaling Tecnologico}

\subsection{Formulazione quantitativa della QSA: tempi, frequenza e lunghezza d'onda}
Per stabilire in modo quantitativo quando sia lecito applicare l'ipotesi di quasi-staticità ed adottare le leggi di Kirchhoff, occorre confrontare la dinamica temporale del segnale con la velocità di propagazione dei campi elettromagnetici nel circuito.

Si considerino le seguenti grandezze fisiche caratteristiche:
\begin{itemize}
    \item $D$: la dimensione spaziale/lineare massima del circuito $[\mathrm{m}]$;
    \item $c$: la velocità di propagazione del campo elettromagnetico nel mezzo (prossima alla velocità della luce nel vuoto, $c \approx 3 \cdot 10^8\,\mathrm{m/s}$) $[\mathrm{m/s}]$;
    \item $v(t) = V_0 \sin(\omega t)$: uno stimolo sinusoidale applicato al circuito con frequenza operativa massima $f_{\text{max}}$ e periodo $T_{\text{max}} = \frac{1}{f_{\text{max}}}$.
\end{itemize}

Il tempo impiegato dal campo elettromagnetico per propagarsi attraverso l'intero circuito è dato da:
\begin{equation}
    t_{\text{prop}} = \frac{D}{c}
\end{equation}

Affinché il campo elettrico e le sue variazioni si trasmettano istantaneamente in tutti i punti del sistema (garantendo l'assenza di ritardi di fase o effetti propagativi lungo i conduttori), il periodo del segnale $T_{\text{max}}$ deve essere molto maggiore del tempo di propagazione $t_{\text{prop}}$:
\begin{equation}
    T_{\text{max}} \gg t_{\text{prop}} \iff \frac{1}{f_{\text{max}}} \gg \frac{D}{c}
\end{equation}

Esprimendo la condizione in termini di \textbf{frequenza operativa massima} $f_{\text{max}}$, si ricava la prima regola empirica fondamentale:
\begin{equation}
    f_{\text{max}} \ll \frac{c}{D}
\end{equation}
Si noti che con "molto minore/maggiore" si intende un fattore $\approx 10^2$.

In maniera del tutto equivalente, ricordando la relazione fondamentale delle onde $c = \lambda \cdot f$, la condizione può essere espressa confrontando la \textbf{lunghezza d'onda del segnale} $\lambda$ con la dimensione fisica $D$:
\begin{equation}
    \lambda = \frac{c}{f_{\text{max}}} \gg D
\end{equation}

\textit{Conclusione fisica}: L'approssimazione QSA (e di conseguenza l'uso delle leggi di Kirchhoff) è valida se la lunghezza d'onda del segnale è nettamente superiore alle dimensioni fisiche del circuito ($\lambda \gg D$).

\subsection{Sistemi a Parametri Concentrati vs. Distribuiti ed Elettronica RF}
In base al rispetto o meno della condizione $\lambda \gg D$, i circuiti si dividono in due grandi categorie:
\begin{enumerate}
    \item \textbf{Sistemi a Parametri Concentrati (\textit{Lumped Parameter Systems})}: Quando la QSA è verificata ($\lambda \gg D$), le tensioni e le correnti dipendono unicamente dal tempo $t$ e sono uniformi lungo ciascun interconnessione o componente. È il campo d'azione dell'elettronica classica e delle leggi di Kirchhoff.
    \item \textbf{Sistemi a Parametri Distribuiti (\textit{Distributed Parameter Systems})}: Quando le dimensioni $D$ diventano comparabili con la lunghezza d'onda ($\lambda \sim D$), le variabili $V(x,t)$ e $I(x,t)$ variano continuamente nello spazio $x$.
    \begin{itemize}
        \item \textit{Esempio classico}: Il cavo coassiale o le linee di trasmissione, in cui ogni tratto infinitesimo di linea $dx$ presenta resistenza $R_{dx}$, induttanza $L_{dx}$ e capacità $C_{dx}$ distribuite, la cui risoluzione richiede l'equazione dei telegrafisti.
    \end{itemize}
\end{enumerate}

Se le condizioni QSA falliscono, il circuito si comporta come un' \textbf{antenna}: emette o riceve radiazione elettromagnetica e genera correnti indotte trasversali. In questo regime, le leggi di Kirchhoff vanno abbandonate ed il circuito deve essere progettato risolvendo le equazioni di Maxwell complete, rientrando nell'ambito dell' \textbf{Elettronica a Radiofrequenza (\textit{RF Electronics})}.

\subsection{Qualche esempio numerico}
Per comprendere l'impatto pratico della condizione $f_{\text{max}} \ll \frac{c}{ D}$, si confrontano due architetture circuitali:

\begin{itemize}
    \item \textbf{Schede elettroniche a componenti discreti (anni '60--'70)}:
    Avente dimensione tipica $D \approx 30\,\mathrm{cm} = 0.3\,\mathrm{m}$.
    \begin{equation}
        t_{\text{prop}} = \frac{0.3\,\mathrm{m}}{3 \cdot 10^8\,\mathrm{m/s}} = 1\,\mathrm{ns} \implies \frac{c}{D} = 1\,\mathrm{GHz}
    \end{equation}
    Applicando il fattore di sicurezza $10^2$, la massima frequenza di lavoro entro cui valgono le leggi di Kirchhoff è limitata a:
    \begin{equation}
        f_{\text{max}} \ll \frac{1\,\mathrm{GHz}}{100} \implies f_{\text{max}} \approx \text{pochi } \mathrm{MHz}
    \end{equation}
    Oltre questa frequenza, le piste della scheda da $30\,\mathrm{cm}$ si comporterebbero come antenne.
    
    \item \textbf{Circuiti Integrati e Microprocessori (CPU attuali)}:
    Per operare con frequenze di clock nell'ordine dei $\mathrm{GHz}$ ($f_{\text{max}} \sim 1 \div 3\,\mathrm{GHz}$) senza ricadere nella complessa teoria dell'elettronica a radiofrequenza o nei parametri distribuiti, è \textbf{fisiologicamente indispensabile} ridurre la dimensione $D$ a pochi millimetri o centimetri ($D \sim \mathrm{cm}$).
\end{itemize}

\subsection{Lo Shrinking Tecnologico e la Legge di Moore}
La necessità di incrementare la frequenza operativa dei circuiti mantenendo valide le leggi di Kirchhoff ha rappresentato la principale spinta verso la miniaturizzazione dell'elettronica moderna (\textit{dalla microelettronica alla nanoelettronica}).

A differenza dei transistori a giunzione bipolare (BJT), la struttura geometrica del transistore a effetto di campo MOS (\textbf{MOSFET}) presenta un'architettura piana particolarmente semplice:
\begin{itemize}
    \item Un substrato di silicio con due pozzetti drogati (Source e Drain) separati da una distanza definita detta lunghezza di canale ($L$);
    \item Un elettrodo superiore di controllo (\textit{Gate}) separato da un sottile strato di biossido di silicio di spessore $t_{\text{ox}}$.
\end{itemize}

\textit{(Nota: torneremo più nel dettaglio su questo componente e sulla sua tecnologia nella seconda parte del corso)}.\\

La semplicità geometrica del MOSFET permette un'operazione di scala dimensionale nota come \textit{shrinking} tecnologico:
\begin{enumerate}
    \item \textbf{Riduzione del Tempo di Transito e Aumento della Velocità}:
    Gli elettroni o le lacune trasportati nel canale si muovono con una velocità di deriva $v_{\text{drift}} = \mu E$. Riducendo la lunghezza del canale $L$, la distanza che i portatori devono percorrere diminuisce drasticamente, riducendo il tempo di transito $\tau_{\text{transit}} \propto \frac{L^2}{\mu V}$ . Di conseguenza, il transistore diventa intrinsecamente più veloce e in grado di commutare a frequenze elevate ($\mathrm{GHz}$).

    \item \textbf{Aumento della Densità di Integrazione e Legge di Moore}:
    La \textbf{Legge di Moore} stabilisce una regola empirica fondamentale secondo cui il numero di transistori integrabili all'interno di un singolo chip di silicio raddoppia indicativamente ogni $1.5 \div 2$ anni. 
    
    Lo \textit{shrinking} della lunghezza di canale $L$ consente di stipare un numero sempre maggiore di dispositivi per unità di superficie ($\mathrm{cm}^2$), incrementando la potenza di calcolo. Il chip di silicio vero e proprio, di dimensioni pari a pochi millimetri o centimetri quadrati, viene infine alloggiato e protetto all'interno di un contenitore plastico o ceramico (\textit{package}, come il classico \textit{Dual In-Line Package} -- DIP).
\end{enumerate}
    
  
% --- CAPITOLI DA INTEGRARE NELLE PROSSIME LEZIONI ---
\chapter{Bipoli Fondamentali e Analisi nel Dominio di Laplace}


% ==========================================
% PARTE II: TECNOLOGIA CMOS
% ==========================================
\part{Tecnologia CMOS}




% ==========================================
% PARTE III: PROGETTAZIONE DI CIRCUITI INTEGRATI
% ==========================================
\part{Progettazione di Circuiti Integrati CMOS}



\end{document}